

\begin{figure}[t!]
\centering
\includegraphics[width=\textwidth]{figs/Fig_experiment-combined.pdf}
\caption{Synthetic experiments.
\textbf{(a)}~Phase behavior as a function of inverse temperature $\beta$ ($d=64$, $K=16$). Left axis (blue): mean cosine similarity to the nearest stored pattern; right axis (coral): scaled entropy $H(\mathbf{a})/\log K$. Both diagnostics reveal a smooth transition centered near $\beta\approx 5\text{--}10$ (gold band).
\textbf{(b)}~Convergence validation ($d=8$, $K=4$, $\beta=5$). Pooled energy density from eight independent chains (gray) overlaid on a long-run reference distribution (coral); inset reports the Kolmogorov--Smirnov statistic and moment differences.}
\label{fig:experiments}
\end{figure}

% ---------- Experiment 1,2 and 3: Temperature Spectrum, Convergence and Load Ratio ----------
\paragraph{Temperature spectrum, Convergence diagnostics and Load Ratio using synthetic data.}
For the first three experiments we constructed memory matrices whose columns are independent draws from the uniform distribution on the unit sphere $\mathbb{S}^{d-1}$.
Concretely, each pattern was sampled as $\mathbf{x}_{k}\sim\mathcal{N}(\mathbf{0},\mathbf{I}_{d})$ and normalized, $\mathbf{x}_{k}\leftarrow\mathbf{x}_{k}/\lVert\mathbf{x}_{k}\rVert$.
This is the standard generative model in the Hopfield capacity literature~\cite{amitStoringInfiniteNumbers1985,folliMaximumStorageCapacity2017} and ensures that, in expectation, the $K$ patterns are nearly orthogonal for moderate load ratios $K/d$.
Unless otherwise noted we fixed $d=64$; the convergence experiment used $d=8$ so that the Boltzmann target density remained tractable for ground-truth comparison.
All datasets were generated with fixed random seeds.

We evaluated how the inverse temperature $\beta$ controlled the transition from noise-dominated diffusion to near-deterministic retrieval (Fig.~\ref{fig:experiments}a).
We stored $K=16$ unit-norm patterns in $\mathbf{X}\in\mathbb{R}^{64\times 16}$ and ran Algorithm~1 for $T=10{,}000$ iterations at each of 25 values of $\beta$ spanning the range $[10^{-2},\,10^{3}]$ (approximately log-spaced), with step size $\alpha=0.01$ and a burn-in of $2{,}000$ samples.
For every post-burn-in sample $\boldsymbol{\xi}_t$ we recorded two diagnostics: the cosine similarity to the nearest stored pattern, $\max_{k}\mathbf{m}_{k}^{\top}\boldsymbol{\xi}_{t}/(\lVert\mathbf{m}_{k}\rVert\lVert\boldsymbol{\xi}_{t}\rVert)$, and the Shannon entropy of the attention weights, $H(\mathbf{p})=-\sum_{k}p_{k}\log p_{k}$ where $\mathbf{p}=\operatorname{softmax}(\beta\mathbf{X}^{\top}\boldsymbol{\xi}_{t})$.

Both curves exhibited a smooth sigmoidal transition centered near $\beta\approx 5\text{--}10$.
At small $\beta$ the cosine similarity plateaued around $0.22$, consistent with the expected inner product between a random unit vector and 16 stored patterns in 64 dimensions, while the entropy saturated at $\log K\approx 2.77$~nats, the uniform limit.
As $\beta$ increased through the transition region the attention weights concentrated on a single pattern: the entropy dropped sharply to zero by $\beta\approx 25$, and the cosine similarity rose to $0.969$ at $\beta=1000$.
The variance of both diagnostics peaked in the transition region, where the chain intermittently switched between pattern-aligned and exploratory states.
These results confirmed that $\beta$ acts as a continuous order parameter interpolating between disordered and ordered phases, in direct analogy to the paramagnetic-to-ferromagnetic transition in classical Hopfield networks~\cite{ackeleyLearningAlgorithmBoltzmann1985}; the useful generation regime occupies intermediate $\beta$ values whose scale is problem-dependent and set by the SNR argument of Eq.~\eqref{eq:snr}.

We verified that Algorithm~1 converged to the correct Boltzmann target by running multiple independent chains on a small system where thorough mixing is feasible (Fig.~\ref{fig:experiments}b).
We stored $K=4$ unit-norm patterns in $\mathbf{X}\in\mathbb{R}^{8\times 4}$ and launched eight chains from different random initializations at $\beta=5$ with step size $\alpha=0.01$ and $T=20{,}000$ iterations, discarding the first $5{,}000$ steps as burn-in.
All eight chains converged to the same stationary energy level; pooled post-burn-in energies closely matched a long-chain reference distribution ($T=200{,}000$, burn-in $50{,}000$) with a small KS statistic, confirming convergence to the correct Boltzmann target rather than a single mode.

We mapped the joint effect of memory load and temperature on attention selectivity by constructing a phase diagram over the load ratio $K/d$ and inverse temperature $\beta$ (Fig.~\ref{fig:phase-diagram}).
We set $d=64$ and swept $K/d$ across 10 values in $[0.05,\,2.0]$ and $\beta$ across 20 log-spaced values in $[0.5,\,100]$, running Algorithm~1 on five independent datasets per condition with $T=10{,}000$ iterations, step size $\alpha=0.01$, and a burn-in of $2{,}000$.
For each $(\beta,\,K/d)$ pair we reported the attention concentration $C = 1 - H(\mathbf{a})/\log K$, averaged over post-burn-in samples and datasets; $C=1$ indicates perfect pattern-selective attention and $C=0$ the uniform limit.
The resulting heatmap revealed a clear phase boundary (dashed $C=0.5$ contour): a retrieval regime in the upper-left corner (high~$\beta$, low~$K/d$) where the sampler locked onto individual patterns, and a diffuse regime in the lower-right where attention spread uniformly across all memories.
The boundary shifted rightward as $\beta$ increased, indicating that higher inverse temperature extends reliable retrieval to larger memory loads, consistent with classical Hopfield capacity theory~\cite{amitStoringInfiniteNumbers1985}.


% ---------- Experiment 4: MNIST Baseline Comparison ----------
\paragraph{Image generation on MNIST.}
The synthetic experiments above validated the sampler's statistical properties on controlled data. We tested whether it could generate structured, recognizable outputs from real-world patterns and compared it against six baselines.
We selected $K=100$ images of the digit ``3'' from the MNIST training set~\cite{lecunGradientBasedLearningApplied1998}, flattened each to a vector in $\mathbb{R}^{784}$, and normalized to unit norm to form the memory matrix $\mathbf{X}\in\mathbb{R}^{784\times 100}$ (load ratio $K/d\approx 0.13$).
Because the energy landscape at high $\beta$ is deeply multimodal, a single long chain can become trapped in one basin for the entire run.
We therefore launched 30 independent chains, each initialized near a different randomly chosen stored pattern with a small Gaussian perturbation ($\sigma_{\mathrm{init}}=0.01$).
Each chain ran Algorithm~\ref{alg:stochastic-attention} at $\beta=2000$ with step size $\alpha=0.01$ for $T=5{,}000$ iterations; after discarding the first $2{,}000$ steps as burn-in, thinning every 100th iteration to obtain 30 post-burn-in snapshots, and sub-sampling 5 evenly spaced snapshots per chain, we retained $30\times 5=150$ generated images that collectively span 30 distinct energy basins.

The choice of $\beta=2000$ follows from the per-step signal-to-noise ratio,
\begin{equation}
\mathrm{SNR}
\;=\;\frac{O(\alpha)}{\sqrt{2\alpha d/\beta}}
\;=\;\sqrt{\frac{\alpha\beta}{2d}}\,,
\label{eq:snr}
\end{equation}
where numerator and denominator are the gradient displacement and expected noise magnitude. The synthetic experiment placed the transition near $\mathrm{SNR}\approx 0.025$ ($\beta\approx5$--$10$, $d{=}64$); inverting at $d{=}784$ gives $\beta\approx 100$, and we set $\beta=2000$ ($\mathrm{SNR}=0.113$, ${\approx}4.5\times$ the transition) to operate well inside the structured regime.

To isolate the contribution of the energy-landscape structure, we compared against five baselines that used the same memory matrix: (i)~\emph{bootstrap resampling}, which draws a stored pattern uniformly at random; (ii)~\emph{Gaussian perturbation}, which selects a random stored pattern and adds isotropic noise $\mathcal{N}(\mathbf{0},\sigma^{2}\mathbf{I})$ with $\sigma$ matched to the sampler's per-step noise scale $\sqrt{2\alpha/\beta}$; (iii)~\emph{random convex combination}, which outputs $\mathbf{X}\mathbf{w}$ for weights $\mathbf{w}\sim\mathrm{Dirichlet}(1,\dots,1)$; (iv)~\emph{GMM-PCA}, which projects the $K$ patterns onto the top-$50$ PCA components, fits a $10$-component diagonal-covariance GMM via EM, and draws samples from the learned distribution (full specification in Appendix~\ref{app:gmm-pca}); (v)~\emph{VAE}~\cite{kingmaWellingAutoEncoding2014} (latent dim 8, encoder 784$\to$256$\to$128$\to\mu/\log\sigma^2$, decoder symmetric), trained for 4{,}000 epochs on the same $K{=}100$ patterns with two-phase training to avoid posterior collapse, samples from $\mathcal{N}(\mathbf{0},\mathbf{I})$ (full specification in Appendix~\ref{app:vae}); and (vi)~\emph{MALA} (Metropolis-Adjusted Langevin Algorithm), which uses exactly the same Langevin proposal as Algorithm~\ref{alg:stochastic-attention} but adds a Metropolis--Hastings accept/reject correction to eliminate discretization bias, run with the identical multi-chain protocol.
For each method we generated 150 samples and evaluated three metrics, novelty~$\mathcal{N}$, diversity~$\bar{\mathcal{D}}$, and mean energy~$\bar{E}$, defined as:
\begin{align*}
\mathcal{N}(\hat{\boldsymbol{\xi}}) = 1 - \max_{k}\cos\!\bigl(\hat{\boldsymbol{\xi}},\,\mathbf{m}_{k}\bigr),\qquad
\bar{\mathcal{D}} = \frac{2}{S(S{-}1)}\!\sum_{i<j}\bigl(1-\cos(\hat{\boldsymbol{\xi}}_i,\hat{\boldsymbol{\xi}}_j)\bigr),\qquad
\bar{E} = \frac{1}{S}\sum_{i=1}^{S} E\!\bigl(\hat{\boldsymbol{\xi}}_i\bigr),
\end{align*}
where $\cos(\mathbf{x},\mathbf{y})=\mathbf{x}^{\top}\mathbf{y}/(\lVert\mathbf{x}\rVert\,\lVert\mathbf{y}\rVert)$ denotes cosine similarity.
Novelty~$\mathcal{N}$ measures departure from the nearest stored pattern (higher is more novel). Diversity~$\bar{\mathcal{D}}$ is the mean pairwise cosine distance among generated samples (higher is more diverse). The mean Hopfield energy~$\bar{E}$ measures proximity to the memory manifold (lower indicates more structured outputs).

\begin{table}[t]
\centering
\caption{Quantitative comparison on MNIST digit ``3'' ($K=100$, $d=784$). $\mathcal{N}$ and $\bar{\mathcal{D}}$ are higher-is-better ($\uparrow$); $\bar{E}$ is lower-is-better ($\downarrow$). Values are mean $\pm$ SE across 30 chains. Two SA rows demonstrate the temperature knob: $\beta{=}2000$ ($\mathrm{SNR}{=}0.113$, structured retrieval) and $\beta{=}200$ ($\mathrm{SNR}{=}0.036$, near transition edge; generation regime). $^\dagger$Energy is positive at $\beta{=}200$ because samples explore off the attractor manifold.}
\label{tab:mnist-baselines}
\small
\begin{tabular}{@{}lccc@{}}
\toprule
Method & $\mathcal{N}$ $\uparrow$ & $\bar{\mathcal{D}}$ $\uparrow$ & $\bar{E}$ $\downarrow$ \\
\midrule
Bootstrap (replay)          & $0.000 \pm 0.000$ & $0.459 \pm 0.011$ & $-0.500 \pm 0.000$ \\
Gaussian perturbation       & $0.004 \pm 0.000$ & $0.450 \pm 0.013$ & $-0.496 \pm 0.000$ \\
Random convex combination   & $0.092 \pm 0.000$ & $0.008 \pm 0.000$ & $-0.399 \pm 0.000$ \\
GMM-PCA ($r{=}50$, $C{=}10$) & $0.198 \pm 0.004$ & $0.419 \pm 0.011$ & $-0.303 \pm 0.005$ \\
VAE (latent${=}8$)          & $0.214 \pm 0.005$ & $0.441 \pm 0.008$ & $-0.286 \pm 0.005$ \\
MALA ($\beta{=}2000$)       & $0.151 \pm 0.001$ & $0.598 \pm 0.001$ & $-0.305 \pm 0.001$ \\
SA ($\beta{=}2000$, retrieval) & $0.152 \pm 0.001$ & $0.600 \pm 0.001$ & $-0.303 \pm 0.001$ \\
\textbf{SA ($\beta{=}200$, generation)} & $\mathbf{0.548 \pm 0.002}$ & $\mathbf{0.885 \pm 0.002}$ & $1.467 \pm 0.008^\dagger$ \\
\bottomrule
\end{tabular}
\end{table}

\begin{figure}[t]
\centering
\includegraphics[width=0.19\textwidth]{figs/Fig_mnist_grid_bootstrap.png}%
\hfill
\includegraphics[width=0.19\textwidth]{figs/Fig_mnist_grid_gaussian.png}%
\hfill
\includegraphics[width=0.19\textwidth]{figs/Fig_mnist_grid_convex.png}%
\hfill
\includegraphics[width=0.19\textwidth]{figs/Fig_mnist_grid_mala.png}%
\hfill
\includegraphics[width=0.19\textwidth]{figs/Fig_mnist_grid_sa.png}

\small
\makebox[0.19\textwidth]{\centering (a) Bootstrap}%
\hfill
\makebox[0.19\textwidth]{\centering (b) Gaussian}%
\hfill
\makebox[0.19\textwidth]{\centering (c) Convex}%
\hfill
\makebox[0.19\textwidth]{\centering (d) MALA}%
\hfill
\makebox[0.19\textwidth]{\centering (e) Ours}
\caption{Generated MNIST digit ``3'' samples ($4\times 4$ grids) at $\beta{=}2000$ ($\mathrm{SNR}{=}0.113$, structured-retrieval regime). Bootstrap outputs are exact copies of stored images. Gaussian perturbation adds unstructured noise. Random convex combinations produce blurry averages. MALA and our ULA-based stochastic attention sampler ($\beta$ controls the operating mode: $\beta{=}2000$ for structured retrieval, $\beta{=}200$ for generation) produce visually indistinguishable diverse, structured digits, confirming that the Metropolis correction is unnecessary at step size $\alpha{=}0.01$.}
\label{fig:mnist-baselines}
\end{figure}

Table~\ref{tab:mnist-baselines} and Figure~\ref{fig:mnist-baselines} summarized the results. Each non-Langevin baseline failed on at least one axis: bootstrap by construction ($\mathcal{N}{=}0$); Gaussian perturbation by negligible noise ($\sigma{\approx}0.003$, indistinguishable from bootstrap); and random convex combination by mode collapse ($\bar{\mathcal{D}}{=}0.008$). Among static baselines, the VAE (latent dim 8, trained on the same $K{=}100$ patterns) is the strongest: $\mathcal{N}{=}0.214\pm0.005$, $\bar{\mathcal{D}}{=}0.441\pm0.008$, surpassing GMM-PCA on both axes. At $\beta{=}200$ (SNR$=$0.036, the generation regime), SA dominates every baseline, including the best learned model (VAE), on both generation metrics simultaneously: novelty $0.548\pm0.002$ ($2.6{\times}$ VAE, $2.8{\times}$ GMM-PCA), diversity $0.885\pm0.002$ ($2.0{\times}$ VAE, $2.1{\times}$ GMM-PCA); energy rises above zero as samples explore off the attractor manifold (see~$^\dagger$). MALA's 99.2\% acceptance rate at both $\beta$ values confirms negligible ULA bias across regimes. The SNR transition band ($0.02$--$0.03$) thus delimits two qualitatively distinct regimes in the same algorithm: \emph{structured retrieval} ($\beta{=}2000$, crisp digits near stored patterns) and \emph{genuine generation} ($\beta{=}200$, blurry-but-recognizable novel digits). To verify that the high diversity at $\beta{=}200$ reflects genuine basin-crossing rather than multi-chain initialization, we ran a single long chain ($T{=}50{,}000$) from a fixed seed; single-chain diversity reached $0.796$, exceeding the 30-chain $\beta{=}2000$ value of $0.600$ (Table~\ref{tab:single-chain}, Appendix~\ref{app:single-chain}).

Appendix~\ref{app:multidigit} confirmed the same ranking on digits ``1'' and ``8'' (SA: $\mathcal{N}{\approx}0.152$--$0.153$, $\bar{\mathcal{D}}{\approx}0.557$--$0.591$; VAE: $\mathcal{N}{\approx}0.180$--$0.209$, $\bar{\mathcal{D}}{\approx}0.397$--$0.408$; MALA acceptance 99.2\%): the VAE remains the strongest non-Langevin baseline on novelty but is far below both Langevin methods on diversity across all digit classes. Appendix~\ref{app:finance} used S\&P~500 log-returns ($d{=}424$, $K{=}2{,}766$) to characterize the novelty--fidelity trade-off on real data: SA achieved novelty $0.768\pm0.001$ (bootstrap: $0.000$) while bootstrap led on marginal fidelity ($99.3\%$ vs.\ $64.2\%$ KS pass); volatility clustering was not reproduced, a principled consequence of fixed-$\beta$ equilibrium sampling. Appendix~\ref{app:simpsons} scales to face images at $d{=}4{,}096$: the SNR rule prescribes $\beta{=}10{,}000$; SA achieves novelty $0.159\pm0.000$ and diversity $0.293\pm0.001$, with ranking identical to MNIST across a $5.2{\times}$ dimension increase.